% Section V -- Over-training reveals anticipated inference demand (module m3_wedge; theory from m7)
\section{Over-Training Reveals Anticipated Inference Demand}\label{sec:wedge}

Section~\ref{sec:tech} estimated the technology from designed variation. This section turns the estimates around and
uses them to read the choices of the labs that release models. A lab that expects its model to process $T$ tokens at inference over its lifetime, and that minimizes lifetime compute
$6ND+2NT$, equates the ratio of the output elasticities of parameters and data to the ratio of lifetime to training
compute (Proposition~\ref{prop:wedge}):
\begin{equation}\label{eq:sV-foc}
w\equiv\frac{\varepsilon_N}{\varepsilon_D}=1+\frac{T}{3D}=\frac{6ND+2NT}{6ND}.
\end{equation}
Two transformations of $w$ recur below: $T/D=3(w-1)$ is planned lifetime inference tokens per training token, and
$w-1=2NT/(6ND)$ is planned lifetime inference compute as a multiple of training compute. Read forward,
equation~\eqref{eq:sV-foc} is the inference-aware allocation rule of \citet{sardana2024beyond}. Read backward, it is a
measurement device: given a technology that delivers the output elasticities at a released model's inputs, the
allocation reveals how much inference its developer planned for, without any data on usage. The logic is that of the
production approach to markups \citep{deloecker2012markups}, which recovers a wedge from the gap between an input's
output elasticity and its cost share. The economics differ: $w$ is the shadow cost of a second use of the same input,
since every parameter is paid for again on each inference token, not a measure of market power. We measure $w$ for
open-weight base models, check it against stated intentions, usage, and rival explanations, and aggregate it.

\subsection{Measurement}\label{sec:wedge-measure}

Under the Chinchilla form, the wedge depends on the technology only through the exponents and the level of the
expansion path:
\begin{equation}\label{eq:sV-logw}
\ln\hat w=(\alpha+\beta)\ln G-\alpha\ln N+\beta\ln D=\frac{\alpha+\beta}{2}\,\ln\frac{M}{M^*(C)},
\end{equation}
where $M=D/N$ is tokens per parameter, $M^*(C)$ is the training-optimal ratio at the model's own compute, and
$(\alpha+\beta)/2=1/\sigma^*-1$ (Online Appendix A, Corollary~\ref{cor:suff}). Three properties make the inversion
usable across labs. First, $w$ does not depend on Hicks-neutral productivity $\omega$, the irreducible loss $E$,
evaluation noise, or the outer exponent $\kappa$ (Proposition~\ref{prop:wedge}): a lab with a better recipe or a
different validation corpus has the same wedge at the same $(N,D)$, and the rejection of $\kappa=1$ in
Section~\ref{sec:tech} matters only through the inner exponents. Second, the level of $\hat w$ is governed by $M^*(C)$, the weakly
identified object of Section~\ref{sec:ident}, whereas curvature only scales the log distance from the path. Third, at a
given compute every technology in this family ranks models by $M$, so ranks should be far more robust than levels.

We evaluate equation~\eqref{eq:sV-logw} under several estimated technologies. The reference is our Huber refit of the
Chinchilla extraction (Section~\ref{sec:tech}; $\alpha=0.347$, $\beta=0.367$, $\sigma^*=0.737$, $M^*(10^{21})=21.4$),
with 400 pairs-bootstrap draws of $(\alpha,\beta,\ln G)$ that carry the joint sampling uncertainty. The alternatives are
the published estimates of \citet{besiroglu2024chinchilla}, Approach~3 of \citet{hoffmann2022training}, the Farseer
sweep \citep{li2025predictableb} with parameter counts including and excluding embeddings, and the RefinedWeb sweep of
\citet{gadre2024language}. For Meta and AI2 models we add the developers' own technologies: Meta's Llama~3 IsoFLOP law
\citep{grattafiori2024llama}, the primal technology fitted to the same IsoFLOP points, and AI2's OLMo ladder
\citep{bhagia2024establishing}. The six common technologies agree on curvature ($\sigma^*$ between 0.72 and 0.80) but
not on the location of the path: $M^*(10^{21})$ ranges from 3.4 (Gadre et al.) to 34 (Hoffmann et al.). For each model
we report the reference $\hat w$ with its 95 percent bootstrap interval and a \emph{technology-sensitivity band} (technology band for short), the union of the
intervals (or points, for published estimates) across technologies. The band is a sensitivity envelope, not an
identified set: it carries no coverage guarantee, and its bootstrap components reflect sampling error in the sweeps, not
the error from extrapolating beyond them.

The choice data are Sample B (Section~\ref{sec:data}): 173 general-purpose open-weight base models in 68 families,
released 2021--2025, with training-token counts documented by their developers and, where available, exact parameter
counts from the released weights; for mixture-of-experts (MoE) models, $N$ is the number of active parameters. $D$
counts tokens processed, including repetitions, in each model's own tokenizer. Because $\partial\ln\hat
w/\partial\ln D=\beta\approx0.37$, a 20 percent difference in tokenizer compression moves $\hat w$ by about 7 percent.
The lifetime objective also prices an inference FLOP like a training FLOP; if serving costs $p$ times as much per FLOP,
the first-order condition becomes $w=1+pT/(3D)$, so $\hat T$ is in training-cost-equivalent tokens.

Most of these models lie outside the designs that identify the technologies. For 55 percent, $M$ exceeds the largest
ratio in the Chinchilla design (341); for 29 percent, it exceeds Farseer's (1,227 with embeddings counted). Seventy-five
percent use more compute than the largest budget of any design in the band (Chinchilla's $1.3\times10^{22}$ FLOP).
Proposition~\ref{prop:pi} quantifies the consequence: at a model the size of Llama~3 8B, uncertainty in $M^*$ accounts
for most of the sampling variance of $\ln\hat w$, and uncertainty in curvature for little of it
(Section~\ref{sec:ident-pi}). We therefore report levels of $\hat T$ as ranges and rest our conclusions on signs and ranks.

\subsection{Results}\label{sec:wedge-results}

\begin{figure}[tp]
\centering
\includegraphics[width=\textwidth]{m3_wedge_fig4}
\caption{Revealed Wedges and Tokens per Parameter}\label{fig:wedge}
\begin{figurenotes}
Revealed wedge $\hat w=\varepsilon_N/\varepsilon_D$ (log scale) against tokens per parameter $M=D/N$ for open-weight
base models. Blue dots: Sample B under the reference technology, with technology bands (vertical gray bars); hollow
squares: mixture-of-experts models (active $N$); orange: family flagships; green diamonds: Meta models under Meta's own
Llama~3 law, linked to their reference wedges; light gray: other open-weight models in \citet{epochai2026data} (token
counts not verified). Dashed line: the reference technology at $C=10^{24}$ FLOP, $\hat w=(M/M^*(C))^{(\alpha+\beta)/2}$.
Horizontal bars: ranges of $M$ in the Chinchilla, \citet{gadre2024language} and Farseer designs. Horizontal line: $w=1$,
the training-optimal allocation.
\end{figurenotes}
\end{figure}

Figure~\ref{fig:wedge} plots $\hat w$ against $M$, and Table~\ref{tab:wedge} reports selected models. Under the
reference technology, the median model in Sample B has $\hat w=3.19$: it was trained as if its developer expected 6.6
lifetime inference tokens per training token, that is, lifetime inference compute of 2.2 times its training compute.
The sign is robust. Ninety-five percent of models have $\hat w>1$, and 75 percent have $\hat w>2$. The lower bound of
the bootstrap interval exceeds 1 for 87 percent of models (80 percent with the budget-clustered bootstrap of
Section~\ref{sec:tech}), and the whole technology band lies above 1 for 79 percent (136 of 173; median band [1.77,
6.92]). In cost terms, the median model used 2.4 times the training compute that would have minimized training cost at
its loss (Farrell cost efficiency 0.41; Section~\ref{sec:framework}).

% Table 6 -- revealed inference demand, selected models (condensed from output/tables/m3_wedge_table5.tex and
% m3_wedge_models.csv; reference technology = chin, band = band_lo/band_hi, own law = w_meta_a2 / w_olmo)
\begin{table}[tp]
\centering
\caption{Revealed Inference Demand for Selected Models}
\label{tab:wedge}
\footnotesize
\setlength{\tabcolsep}{3.2pt}
\begin{tabular}{@{}lrrrcrrcc@{}}
\toprule
 & $N$ & $D$ & & $\hat w$ & & & Technology & Own \\
Model & (B) & (T) & $M$ & [95\% CI] & $\hat T/D$ & CE & band & law \\
\midrule
\multicolumn{9}{@{}l}{\textit{Panel A. Meta}} \\
Llama 2 70B$^{f}$ & 69.0 & 2.00 & 29 & 1.19 [0.93, 1.67] & 0.6 & 0.98 & [0.55, 3.77] & 0.98 \\
Llama 3 8B & 8.03 & 15.0 & 1,868$^{e}$ & 5.27 [3.97, 8.12] & 12.8 & 0.18 & [2.48, 12.0] & 3.09 \\
Llama 3 70B & 70.6 & 15.0 & 213 & 2.48 [1.81, 3.81] & 4.4 & 0.58 & [1.02, 7.55] & 1.63 \\
Llama 3.1 405B$^{f}$ & 406 & 15.6 & 38 & 1.37 [0.99, 2.22] & 1.1 & 0.93 & [0.42, 5.73] & 0.98 \\
\addlinespace
\multicolumn{9}{@{}l}{\textit{Panel B. Other open-weight families}} \\
Qwen2.5 0.5B & 0.494 & 18.0 & 36,435$^{e}$ & 14.8 [11.4, 22.7] & 41.5 & 0.02 & [6.20, 23.7] & \\
Qwen2.5 72B$^{f}$ & 72.7 & 18.0 & 248 & 2.62 [1.91, 4.07] & 4.9 & 0.54 & [1.06, 8.01] & \\
Qwen3 0.6B & 0.596 & 36.0 & 60,398$^{e}$ & 17.9 [13.5, 28.5] & 50.8 & 0.01 & [6.95, 29.4] & \\
Qwen3 32B$^{f}$ & 32.8 & 36.0 & 1,099$^{c}$ & 4.46 [3.22, 7.33] & 10.4 & 0.24 & [1.98, 11.9] & \\
Gemma 3 1B$^{\dagger}$ & 1.00 & 2.00 & 2,000$^{e}$ & 5.19 [4.24, 7.00] & 12.6 & 0.18 & [3.11, 8.98] & \\
Gemma 3 27B$^{f\dagger}$ & 27.0 & 14.0 & 518$^{c}$ & 3.37 [2.52, 5.20] & 7.1 & 0.38 & [1.62, 8.98] & \\
OLMo 2 1B & 1.48 & 4.05 & 2,727$^{e}$ & 5.86 [4.67, 8.30] & 14.6 & 0.15 & [1.01, 10.6] & 2.01 \\
OLMo 2 32B$^{f}$ & 32.2 & 6.15 & 191 & 2.34 [1.78, 3.42] & 4.0 & 0.61 & [0.23, 6.42] & 0.91 \\
SmolLM2 135M & 0.135 & 2.00 & 14,868$^{e}$ & 10.4 [8.71, 13.7] & 28.2 & 0.04 & [4.93, 14.7] & \\
DeepSeek-V3$^{\ddagger}$ & 37.0 & 14.8 & 400$^{c}$ & 3.09 [2.29, 4.73] & 6.3 & 0.43 & [1.40, 8.54] & \\
\addlinespace
\multicolumn{9}{@{}l}{\textit{Panel C. Stated compute-optimal, Kaplan-era, and fixed-$D$ designs}} \\
Chinchilla 70B & 70.0 & 1.40 & 20 & 1.04 [0.82, 1.42] & 0.1 & 1.00 & [0.49, 3.28] & \\
GPT-3 175B & 175 & 0.300 & 1.7 & 0.43 [0.35, 0.57] & $-$1.7 & 0.61 & [0.19, 1.66] & \\
Pythia 12B$^{f}$ & 11.8 & 0.300 & 25 & 1.09 [0.93, 1.36] & 0.3 & 0.99 & [0.75, 2.66] & \\
\bottomrule
\end{tabular}

\medskip
\begin{tablenotes}
$N$ is the parameter count in billions (active parameters for the mixture-of-experts model$^{\ddagger}$); $D$ is
training tokens in trillions; $M=D/N$. $\hat w=\varepsilon_N/\varepsilon_D$ at the model's $(N,D)$ under the reference
technology (Huber refit of the Chinchilla extraction; Section~\ref{sec:tech}), with a 95 percent pairs-bootstrap interval
(400 draws). $\hat T/D=3(\hat w-1)$ is planned lifetime inference tokens per training token, and $\hat w-1$ is planned
lifetime inference compute as a multiple of training compute; negative values are not inference demand. CE
$=C_{\min}(L)/(6ND)$ is Farrell cost efficiency relative to the training-only frontier. Technology band: union of the
95 percent intervals (or point estimates) for $\hat w$ under the Chinchilla refit, \citet{besiroglu2024chinchilla},
\citet{hoffmann2022training} (Approach 3), Farseer with $N$ including and excluding embeddings, and \citet{gadre2024language}
(RefinedWeb), plus the developer's own technologies for Meta and AI2 models; it is a sensitivity envelope, not a
confidence set. Own law: $\hat w$ under Meta's Llama~3 IsoFLOP law (Approach 2 with $\kappa=1$) for Meta models and under
AI2's OLMo ladder for OLMo~2. $^{f}$Flagship (largest model of its family). $^{\dagger}$Distilled from a larger model.
$^{c}$/$^{e}$: $M$ above the largest ratio in the Chinchilla design (341) / also above that in the Farseer design with
embeddings counted (1,227). Llama~3 and Llama~3.1 8B and 70B have identical $N$ and $D$.
\end{tablenotes}
\begin{tablenotes}[Source]
Panels A and B: Sample B (Section~\ref{sec:data}; parameter counts from the released weights, token counts from
developer documentation). Panel C: \citet{epochai2026data} for Chinchilla and GPT-3 (GPT-3 tokens from
\citealt{brown2020language}); \citet{biderman2023pythia} for Pythia. Authors' calculations.
\end{tablenotes}
\end{table}


Levels depend on the technology. The median $\hat w$ ranges from 1.91 (Hoffmann et al.) and 2.12 (Meta's law) through
2.41--2.67 (Farseer) and 3.16 (Besiroglu et al.) to 4.03 (Gadre et al.), and freeing the outer exponent on Farseer gives
3.85. On the Chinchilla data alone, nonlinear least squares in levels gives 4.11, and keeping the five outlying runs gives
4.78. These differences dwarf sampling error: the
median bootstrap interval spans only 0.80 to 1.39 times the point estimate. As Proposition~\ref{prop:pi} predicts, the
location of the expansion path, not its curvature, drives the uncertainty. Ranks barely move: the Spearman correlation
between the reference wedges and those under each alternative technology lies between 0.95 and 1.00. The tight curve in
Figure~\ref{fig:wedge} has the same source. Under the reference technology $M^*(C)$ changes little with compute (21.4 at
$10^{21}$ and 17.6 at $10^{24}$ FLOP), so $\hat w$ is nearly a function of $M$ alone.

Table~\ref{tab:wedge} makes the magnitudes concrete. Llama~3 8B was trained on 15 trillion tokens ($M=1{,}868$) and has
$\hat w=5.27$ [3.97, 8.12]: its allocation rationalizes about 13 lifetime inference tokens per training token ($\hat
T/D=12.8$ [8.9, 21.4]; $\hat T\approx1.9\times10^{14}$ tokens). Under Meta's own law, which imposes $\kappa=1$ on
Meta's Approach-2 estimates, $\hat w=3.09$ and $\hat T/D=6.3$; the primal fitted to the same IsoFLOP points gives
$\hat w=4.82$, and the technology band for $\hat T/D$ is [4.4, 33]. The model's cost efficiency is 0.18: Meta spent 5.6
times the training compute needed for the same loss, which is rational only if the model was expected to serve many
tokens. Meta's flagships tell the opposite story. Llama~3.1 405B has $\hat w=0.98$ under Meta's law (1.37 [0.99, 2.22]
under the reference technology), as does Llama~2 70B (0.98; reference 1.19 [0.93, 1.67]). Meta trained its largest
models on its own training-optimal path and over-trained the small ones. At the extreme, Qwen3 0.6B saw 36 trillion
tokens ($M=60{,}398$; $\hat w=17.9$), more than 20 times the largest ratio in any design; its level is not credible, but
its rank is.

Where does extrapolation push the levels? \citet{sardana2024beyond}, who train at up to about 10,000 tokens per
parameter, find that fitted Chinchilla laws overstate the loss reduction from extra tokens at extreme ratios. If so, the
true $\varepsilon_D$ at high $M$ is smaller than the extrapolated one, and $\hat w$ understates $w$. On Farseer, the one
near-factorial public sweep, the nonparametric local elasticity of substitution falls with $M$ at given $N$, which
suggestively points the same way. Shape restrictions alone deliver a lower bound on $T$ (Proposition~\ref{prop:pi}), but
on current support it is uninformative: the design points with more parameters and fewer tokens than a Llama-scale model
are themselves under-trained ($\hat w\le0.77$), so the bound implies only $T\ge0$.
\textcolor{red}{[TBD-m9: In our controlled sweep, which reaches $M\approx$~X, the Chinchilla form fitted on $M\le$~Y
over/understates the marginal value of tokens at $M\approx1{,}000$--$2{,}000$ by Z, so reference wedges in this range are
under/overstated by about W percent.]}

\subsection{Stated Intent}\label{sec:wedge-intent}

The wedge reads developers' stated objectives correctly. Chinchilla, designed to be training-compute-optimal, has $\hat
w=1.04$ [0.82, 1.42]. The seven Cerebras-GPT models, trained at the 20 tokens per parameter that their developers call
compute-optimal \citep{dey2023cerebras}, have 0.92--1.00. LLaMA-1 7B and 13B, designed for the best performance at given
inference budgets \citep{touvron2023llama}, have 1.65--2.07, and Llama~3 8B and 70B, trained well beyond the
compute-optimal token count to perform better at inference \citep{grattafiori2024llama}, have 2.48--5.27. Pythia is a
useful warning: every model in the suite saw the same 300 billion tokens by research design \citep{biderman2023pythia},
and the implied wedges range from 1.09 (12B) to 6.49 (70M). Over-training can be the by-product of a fixed-$D$ design
rather than of demand. These five cases were chosen because the rule is stated; they are a sanity check without
statistical power.

\subsection{Trends}\label{sec:wedge-trends}

\begin{figure}[tp]
\centering
\includegraphics[width=\textwidth]{m3_wedge_trends}
\caption{The Revealed Wedge over Time}\label{fig:trends}
\begin{figurenotes}
Production-scale language models ($6ND\ge10^{21}$ FLOP) released 2019--2026 ($n=335$): Sample B and open and closed
models from \citet{epochai2026data}, under the reference technology. Panel (a): median $\hat w$ (log scale) and
interquartile range by release year, for open-weight and closed or API-only models. Panel (b): shares of models with
$\hat w<1$ and $\hat w>2$; the number of models per year is shown at the bottom.
\end{figurenotes}
\end{figure}

Figure~\ref{fig:trends} turns to time, in a production-scale universe of 335 language models with $6ND\ge10^{21}$ FLOP
released 2019--2026: Sample B plus open and closed models from Epoch AI's database whose $N$ and $D$ are recoverable
\citep{epochai2026data}, 271 open-weight and 64 closed. The industry moved from under- to over-training within three
years. The share of models with $\hat w<1$ was 47 percent in 2021 and 41
percent in 2022, fell to 18 percent in 2023, and was 0 percent in 2024 and 2 percent in 2025. The median $\hat w$ rose
from 1.06 in 2022 to 1.83 in 2023, 3.45 in 2024, and 4.68 in 2025, as median tokens per parameter rose from 23 to 1,303.
The early under-training is concentrated in closed giants: GPT-3 175B (0.43), Gopher (0.37), Megatron-Turing NLG 530B
(0.28), and PaLM 540B (0.41). It is not negative demand. By Proposition~\ref{prop:wedge}, a wedge below one can reveal
only a binding data constraint or optimization against a mistaken technology, here the law of
\citet{kaplan2020scaling}, which prescribed far fewer tokens per parameter.

Open-weight developers led the shift. Conditional on release year and MoE architecture, open-weight models have $\ln\hat
w$ higher by 0.57 (standard error 0.09), and by 0.44 (0.09) when $\ln C$ is also held fixed (standard errors clustered by
developer, 94 clusters). Dropping 11 Epoch rows whose token counts are inconsistent with Epoch's own compute estimates
gives 0.59 and 0.43. Among developers with at least four production-scale models, the median $\hat w$ is 6.6 at Hugging
Face, 4.8 at Alibaba, 2.3 at Google, 1.8 at Meta, and 1.0 at Cerebras. Two selection caveats apply. Closed developers
rarely disclose $N$ and $D$ (64 closed models, 3 of them from 2025), and some open developers do not disclose $D$
(Mistral's models are absent), so the premium compares disclosed allocations.

\subsection{Within-Family Revealed Preference and Over-Identification}\label{sec:wedge-family}

Families offer a way around $M^*$, the least reliable input to the inversion. In the spirit of
\citet{gandhi2020identification}, a first-order condition at one point can calibrate the technology. If a family's
flagship, its largest model, sits on its developer's training-optimal path ($w_f=1$), each sibling's wedge relative to
the flagship is
\begin{equation}\label{eq:sV-rel}
w_i^{\mathrm{rel}}=\Big(\frac{N_i}{N_f}\Big)^{-\alpha}\Big(\frac{D_i}{D_f}\Big)^{\beta}.
\end{equation}
This ratio depends on the exponents but not on the family's $A/B$, so it removes the dominant source of uncertainty in
Proposition~\ref{prop:pi}. Smaller siblings are more over-trained than their flagships with few exceptions: 95 percent of
105 non-flagship siblings in 44 families have $w^{\mathrm{rel}}>1$; the exceptions are five small Cerebras-GPT models
trained at a constant $M=20$ (0.92--0.98). Given the exponents and the observed token counts, this is close to mechanical.

The normalization is itself testable, and after 2024 it fails. The 18 flagships released before 2024 have a median
reference $\hat w$ of 1.14, and half of their bootstrap intervals contain 1. The 26 flagships released from 2024 onward
have a median of 3.14, and 92 percent of their intervals lie above 1 (Qwen3 32B, for example, has 4.46 [3.22, 7.33]).
Recent flagships are themselves over-trained, so flagship-normalized demands are lower bounds. Meta is the exception:
under its own law, both Llama~2 70B and Llama~3.1 405B have $\hat w=0.98$.

Families also permit an over-identification test in the spirit of \citet{raval2023testing}. If every family faced one
iso-elastic demand for model size, $T_i=\tau_fN_i^{-\eta}$, the slope of $\ln\hat T_i$ on $\ln N_i$ would be common
across families. The pooled slope is $\hat\eta=0.53$ (0.07), and homogeneity is rejected ($F=2.55$, $p=0.001$; 138
models in 42 families; a classical $F$ test on models with $\hat T>0$). Much of this rejection is mechanical, because a
family that trains all sizes on the same $D$ has $\hat T_i/D$ a nonlinear function of $N_i$ alone. In the 18 families
whose token counts vary by more than 10 percent across sizes, the rejection is marginal ($p=0.047$ with absolute $\hat
T$, 58 models; $p=0.12$ with flagship-normalized $\hat T$, 33 models in 11 families). Families that scale $D$ with $N$
(e.g., Gemma~2, SmolLM2, Phi-3.5) imply planned demand rising with size: either their large models were expected to
serve more tokens, or these developers set $D$ by per-size budget rules rather than by minimizing lifetime cost. The
test cannot distinguish the two, and it does not reject the wedge itself.

\subsection{Validation against Usage}\label{sec:wedge-validation}

If over-training reflects anticipated demand, more over-trained models should be used more. Because $w$ is a function of
$(N,D)$, the informative comparison holds compute fixed: at a given $C$, a higher $M$ means a smaller model, cheaper to
serve but, under a common technology, with higher loss than the compute-optimal one. Table~\ref{tab:validation}
regresses usage on $\ln M$, controlling for $\ln C$, release year, and log months since release, for 164 Sample B models
from 26 developers. Downloads rise roughly one for one with tokens per parameter: the coefficient is 0.97 (0.18) for
all-time and 1.13 (0.19) for 30-day downloads, and 0.72 (0.14) and 0.94 (0.14) with developer fixed effects. The number
of OpenRouter providers serving a model rises by 0.08 (0.04) log points per log point of $M$, and LMArena votes by 0.29
(0.16) among 46 matched models. Without controls, the correlations with likes (0.13, standard error 0.25) and votes
($-0.17$, 0.40) are weak or negative: the size confound at work.

% Table 7 -- usage validation (from output/tables/m3_wedge_validation.csv; specs uncond, cond_C, cond_C_lab)
\begin{table}[tp]
\centering
\caption{Over-Training and Usage}
\label{tab:validation}
\footnotesize
\setlength{\tabcolsep}{4pt}
\begin{tabular}{@{}lcccccc@{}}
\toprule
 & All-time & 30-day & & Listed on & OpenRouter & LMArena \\
 & downloads & downloads & Likes & OpenRouter & providers & votes \\
 & (1) & (2) & (3) & (4) & (5) & (6) \\
\midrule
\multicolumn{7}{@{}l}{\textit{Panel A. No controls}} \\
$\ln\hat w$ & 1.54 & 2.60 & 0.13 & 0.112 & 0.136 & $-$0.17 \\
 & (0.40) & (0.34) & (0.25) & (0.040) & (0.047) & (0.40) \\
\addlinespace
\multicolumn{7}{@{}l}{\textit{Panel B. Controls: $\ln C$, release year, log age}} \\
$\ln M$ & 0.97 & 1.13 & 0.27 & 0.055 & 0.081 & 0.29 \\
 & (0.18) & (0.19) & (0.16) & (0.030) & (0.036) & (0.16) \\
Models & 164 & 164 & 164 & 164 & 164 & 46 \\
Developers & 26 & 26 & 26 & 26 & 26 & 12 \\
\addlinespace
\multicolumn{7}{@{}l}{\textit{Panel C. Panel B and developer fixed effects}} \\
$\ln M$ & 0.72 & 0.94 & 0.29 & 0.042 & 0.059 & 0.25 \\
 & (0.14) & (0.14) & (0.16) & (0.033) & (0.044) & (0.14) \\
Models & 158 & 158 & 158 & 158 & 158 & 42 \\
Developers & 20 & 20 & 20 & 20 & 20 & 8 \\
\bottomrule
\end{tabular}

\medskip
\begin{tablenotes}
Each cell is the coefficient from a separate OLS regression; standard errors clustered by developer in parentheses.
Sample B models with Hugging Face metadata. Outcomes: $\ln(1+{}$downloads$)$, summed over the base repository and the
developer's official post-trained releases (columns 1--2); $\ln(1+{}$likes$)$ of the base repository (3); an indicator
for being listed on OpenRouter, as a linear probability model (4); $\ln(1+{}$number of distinct OpenRouter
providers$)$ (5); $\ln(1+{}$votes$)$ of the matched release in the LMArena text arena (6). Log age is the log of
months since release. Panel A uses $\ln\hat w$ under the reference technology, which is a function of $(\ln N,\ln D)$
and therefore confounds size. Given $\ln C$, $\ln\hat w=[(\alpha+\beta)/2][\ln M-\ln M^*(C)]$ and
$\ln N=(\ln C-\ln6-\ln M)/2$, so Panels B and C compare models of equal compute that differ in size and wedge.
Panel C drops developers with a single model. In Panel A, column 6 uses the 46 matched models and the other columns
use 164.
\end{tablenotes}
\begin{tablenotes}[Source]
Hugging Face Hub (snapshot 2026-09-24), \citet{openrouter2026models} (snapshot 2026-09-23),
\citet{lmarena2026leaderboard}; authors' calculations.
\end{tablenotes}
\end{table}


This is weak evidence of demand. At fixed $C$, the coefficient on $\ln M$ equivalently says that smaller models of equal
compute are downloaded more, which the larger number of users able to run small models also predicts. Developer or
vintage quality correlated with $M$ is not ruled out; the fixed-effects estimate (0.72) is the conservative number.
Downloads are not inference tokens, and OpenRouter lists only 21 of the 164 models in a current snapshot. We read
Table~\ref{tab:validation} as a correlation consistent with anticipated demand, not as an estimate of $T$; per-model
inference volumes, such as the OpenRouter panel of \citet{demirer2025emerging}, would permit a test of levels.

\subsection{Rival Wedges}\label{sec:wedge-rivals}

Proposition~\ref{prop:wedge} signs the rival explanations. A binding data constraint lowers $w$; a binding memory or
latency constraint raises it with no inference demand at all; and factor-augmenting productivity multiplies the
econometrician's wedge by $e^{-\chi}$, where $\chi=\beta\psi_D-\alpha\psi_N$ is the factor-bias index. Online Appendix Table~\ref{tab:app-rivals}
gives details.

Data scarcity and Kaplan-era beliefs produce $w<1$, so they cannot generate the modern pattern; under-training is a
pre-2024 phenomenon (one of the 171 production-scale models released since 2024 has $\hat w<1$). Hardware tiers are a
real but second-order rival. Of the open-weight production-scale models released 2023--2026, 26 percent have 6.5--9.5
billion parameters, which fit on one 16--24 GB accelerator at 16-bit precision, against 9 percent under a log-normal fit.
Given $\ln C$ and year, these models have $\ln\hat w$ higher by 0.11 (0.06), about a tenth of the median $\ln\hat w$ of
1.16. Under the parameters of \citet{besiroglu2024chinchilla}, a memory cap at half the training-optimal size would by
itself produce $w=1.78$ with $T=0$, so tiers cannot be dismissed, but the estimated effect is too small to explain the
level.

Non-homotheticity is the most serious rival: if the training-optimal ratio rises with compute, large models trained at
high $M$ may be closer to optimal than the Chinchilla form implies. Farseer's own non-homothetic form
\citep[Eq.~3]{li2025predictableb} implies $M^*(C)$ rising from 25 at $10^{20}$ to 82 at $10^{23}$ FLOP (195 at $10^{24}$,
extrapolated). It can rationalize the token counts of large models but not of small ones. Within Farseer's parameter
support (66 models), its wedges are higher than those of the Chinchilla form fitted to the same sweep (median 4.12 versus
3.44); above that support (75 models), they are lower (1.54 versus 2.12), and Llama~3 70B would be under-trained (0.62).
This last comparison extrapolates a nine-parameter form one to two orders of magnitude beyond its support.

Distilled and synthetic-data-heavy models show no detectable difference in $\ln\hat w$ given $\ln C$ and year (0.02
(0.23) and 0.06 (0.18)); the test has no power, and the sign of any bias is ambiguous because distillation could augment
either input. Factor-biased recipes are bounded by the estimates of Section~\ref{sec:tech}: across 25 DataDecide recipes,
factor bias moves the wedge at given $(N,D)$ by a factor of at most 1.25--1.29 (1.80 in the most extreme robustness
variant; 1.02--1.04 across the corpora of Gadre et al.). Eighty-eight percent of Sample B have $\hat w$ above 1.29, and 79
percent above 1.80; for 69 and 49 percent, the whole technology band lies above these factors. Recipe bias of the
measured size cannot make the typical modern model training-optimal, although, combined with the widest band and the
extreme factor, it could for about half of the sample.

\subsection{An Illustrative Aggregate}\label{sec:wedge-aggregate}

The compute-weighted mean of $w-1$,
\begin{equation}\label{eq:sV-agg}
\frac{\sum_i(w_i-1)C_i}{\sum_iC_i}=\frac{\sum_i2N_iT_i}{\sum_i6N_iD_i},
\end{equation}
is the planned lifetime inference compute of a set of models as a multiple of their training compute. For the 271
open-weight production-scale models it is 2.0 [1.2, 3.7] under the reference technology (technology band [0.4, 7.2]).
For 2025 releases it is 3.2 [2.1, 5.7] (band [1.05, 9.6]), up from 1.2 [0.6, 2.3] for 2024. For 2022 releases it is
negative ($-0.40$), because BLOOM, OPT-175B, and GLM-130B were under-trained (0.03 with $\hat T$ truncated at zero).
Three caveats make equation~\eqref{eq:sV-agg} an illustration of the object rather than a measurement: it is dominated
by the largest models, whose wedges are the most extrapolated; it depends heavily on the technology, from 0.67 (Hoffmann
et al.) to 3.26 (Gadre et al.); and we found no public per-model inference volumes against which to validate it.

The robust findings of this section are therefore ordinal. The typical open-weight model is over-trained under every
technology we consider, small models more than large ones and open developers more than closed ones, and the industry
switched from under- to over-training in about three years. The cardinal findings, how many inference tokens a given
developer expected, inherit the fragility of $M^*$ documented in Sections~\ref{sec:ident} and~\ref{sec:tech}; calibrating
$M^*$ within families and experimenting at high $M$ are the natural ways to tighten them.
